% Notas de aula — Capítulo 9: Boletins Meteorológicos e Análise de Cartas
% Curso: Meteorologia Prática (baseado em R. Stull, Practical Meteorology, CC BY-NC-SA 4.0)
% Caixas verdes = conteúdo literal do quadro; linhas verdes = encenação.
\documentclass[11pt]{article}
\usepackage{../comum/preambulo}
\graphicspath{{../figuras/cap09/}}

\title{Capítulo 9 --- Boletins Meteorológicos e Análise de Cartas\\
{\large Notas de aula (roteiro de quadro-negro)}}
\author{Curso de Meteorologia Prática}
\date{}

\begin{document}
\maketitle
\thispagestyle{fancy}

\noindent\textbf{Plano sugerido:} 1 aula de 2\,h $+$ 1 aula prática
(análise manual de uma carta real).
\emph{Aula 1:} \S\ref{sec:metar}--\S\ref{sec:plot} (códigos e o modelo
de estação).
\emph{Aula 2 (prática):} \S\ref{sec:analise} com carta impressa e lápis.

\section{Por que códigos e cartas ainda importam}

\begin{noquadro}[abertura]
\textbf{Cap.~9 --- Boletins e Cartas}\\[2pt]
``Antes de confiar no modelo numérico, o meteorologista confere a
observação — e a observação fala em código.''\\[2pt]
METAR (horário, aviação) $\to$ modelo de estação (um boletim em
2\,cm$^2$) $\to$ carta sinótica (o campo)
\end{noquadro}

Este é o capítulo mais operacional do curso: alfabetização profissional.
A rede mundial de observação — estações, navios, boias, radiossondas,
aviões e os satélites do capítulo anterior \javimos{8}{satélites} —
converge quatro vezes por dia (00/06/12/18\,UTC
\javimos{1}{tempo em UTC}) para as cartas; a mesma torrente de dados
alimenta a assimilação dos modelos \veremos{20}{assimilação de dados}.

\section{Redução ao nível do mar}\label{sec:reducao}

\begin{formadif}[a hipsométrica de novo]
Estações em altitudes diferentes medem pressões incomparáveis (10\,kPa
a cada $\sim$840\,m! \javimos{1}{hidrostática}). Reduz-se tudo ao nível
do mar com a hipsométrica \javimos{1}{equação hipsométrica}:
\[ P_{nm} = P_{est}\,\exp\!\left(\frac{g\,z_{est}}{\Re_d\,\bar T_v}\right), \]
usando a $\bar T_v$ média da coluna fictícia entre a estação e o nível
do mar.
\end{formadif}

\begin{noquadro}[sem redução não há carta]
sem reduzir: mapa de pressão $=$ mapa de \emph{altitude} das estações\\
com redução: as diferenças que sobram são \emph{meteorologia}\\[2pt]
(Caso~4 do notebook mostra o antes/depois com estações do litoral ao
planalto)
\end{noquadro}

O ponto fraco: que $\bar T_v$ usar para uma coluna de ar que não existe?
Numa manhã de inversão \javimos{5}{inversões} a $T$ da estação é fria
demais e a redução \emph{infla} a pressão — o exercício T1 quantifica.

\section{METAR: o código horário da aviação}\label{sec:metar}

\quadro{Decodificar no quadro, grupo a grupo, um METAR real de
Guarulhos. Cada grupo aponta para um capítulo anterior do curso.}

\begin{center}
\texttt{SBGR 221800Z 12008KT 8000 -RA BKN015 OVC080 22/19 Q1016}
\end{center}

\begin{noquadro}[decodificação]
\begin{tabular}{ll}
\texttt{SBGR 221800Z} & estação, dia 22, 18:00 UTC \\
\texttt{12008KT} & vento \emph{de} $120°$, 8\,kt
\javimos{1}{direção do vento} \\
\texttt{8000} & visibilidade 8\,000\,m \\
\texttt{-RA} & chuva fraca ($-$/$+$; RA, DZ, TS, FG\ldots) \\
\texttt{BKN015 OVC080} & 5--7 oktas a 1\,500\,ft; 8 a 8\,000\,ft
\javimos{6}{oktas e teto} \\
\texttt{22/19} & $T$, $T_d$ $\Rightarrow$ UR, LCL
\javimos{4}{ponto de orvalho} \\
\texttt{Q1016} & QNH 1016\,hPa (reduzida ao nível do mar) \\
\end{tabular}
\end{noquadro}

Parser em Python no Caso~1 do notebook — decodificar vira uma linha; o
N1 o estende para rajadas, \texttt{VRB} e \texttt{CAVOK}.

\section{O modelo de estação (station plot)}\label{sec:plot}

\quadro{Desenhar o modelo de estação GIGANTE no quadro (o oficial da
Fig.~9.14 do livro) e preenchê-lo com o METAR decodificado acima. Os
alunos copiam — este desenho é ferramenta de prova.}

\begin{noquadro}[o layout oficial]
centro: círculo de cobertura (oktas)\\
NW: $T$ \quad SW: $T_d$ \quad W: tempo presente (símbolo)\\
NE: pressão codificada (3 dígitos $=$ décimos de hPa sem o milhar)\\
E: tendência de pressão (3\,h) \quad haste: barbela de vento\\[2pt]
barbela: bandeira $=50$\,kt, traço longo $=10$\,kt, curto $=5$\,kt;
a haste aponta \emph{de onde} o vento vem
\end{noquadro}

Decodificação da pressão: 164 $\to$ 1016{,}4\,hPa; 986 $\to$
998{,}6\,hPa — escolher o que dá valor realista (o T3 discute quando a
regra falha: furacões intensos \veremos{16}{pressões extremas}).
O MetPy desenha o modelo oficial (Caso~2 do notebook).

\section{Análise de cartas}\label{sec:analise}

\quadro{Regras de ouro no quadro; depois TODOS analisam a mesma carta
impressa a lápis. Comparar resultados no fim — as diferenças são a
lição.}

\begin{formalivro}[regras de análise]
\begin{enumerate}
  \item Isóbaras a cada 4\,hPa (\ldots 1008, 1012, 1016\ldots), suaves,
        nunca se cruzam;
  \item interpolar entre estações proporcionalmente;
  \item fechar centros: \textbf{A}(lta) e \textbf{B}(aixa) — H/L em
        cartas internacionais;
  \item frentes nas dobras das isóbaras e nos contrastes de
        $T$/$T_d$/vento \veremos{12}{frentes};
  \item isóbaras apertadas $=$ vento forte (o porquê é o
        Cap.~10 \veremos{10}{vento geostrófico}).
\end{enumerate}
\end{formalivro}

Também se analisam cartas de altitude (ex.\ 50\,kPa): ali as isolinhas
são de \emph{altura geopotencial} da superfície isobárica — colinas e
vales cuja inclinação, veremos, é o próprio vento
\veremos{10}{aproximação geostrófica}; os cavados de 50\,kPa comandam os
ciclones de superfície \veremos{13}{ciclogênese}.

À mão isso é treino de olho; em máquina é \emph{análise objetiva} —
interpolar observações irregulares para uma grade. O método clássico de
correções sucessivas (Cressman/Barnes) está no Caso~3 do notebook: é o
ancestral direto da assimilação de dados moderna
\veremos{20}{assimilação}. O N3 acha o raio de influência ótimo por
validação cruzada — o dilema viés$\times$variância ao vivo.

\section{Síntese da aula}

\begin{noquadro}[mapa final --- canto do quadro]
\[ \text{METAR} \to \text{modelo de estação} \to
   \text{redução ao nível do mar} \to \text{isóbaras/análise} \]
próxima aula: \emph{por que} o vento segue as isóbaras (Cap.~10)
\end{noquadro}

\section{Exercícios complementares}\label{sec:exercicios}

Soluções (professor) em \texttt{cap09\_solucoes.ipynb}.

\subsection*{Teóricos}

\begin{enumerate}[label=\textbf{T\arabic*.},itemsep=4pt]
  \item Deduza a fórmula de redução ao nível do mar a partir da
        hipsométrica e calcule a correção para São Paulo ($z = 800$\,m,
        $P_{est} = 92{,}5$\,kPa, $T = 20\,°$C). Discuta o erro de usar a
        $T$ da estação numa manhã de inversão.
  \item Decodifique à mão:
        \texttt{SBSP 221500Z 33012G25KT 9999 TS FEW030CB 28/21 Q1012} e
        desenhe o modelo de estação correspondente (com a barbela
        correta).
  \item A pressão codificada ``052'' num station plot pode ser 1005{,}2
        ou 905{,}2\,hPa. Formule a regra de desempate e dê um exemplo em
        que ela falharia (furacão intenso).
  \item Mostre que, para espaçamento de isóbaras $\Delta n$ na carta, o
        vento geostrófico (prévia do Cap.~10) escala como $1/\Delta n$ —
        e por isso ``apertado $=$ ventoso''.
  \item No método de Barnes, o peso é $w = e^{-d^2/R^2}$. Mostre que
        $R$ grande filtra escalas pequenas (média quase uniforme) e $R$
        pequeno superajusta às observações; relacione com o dilema
        viés$\times$variância.
\end{enumerate}

\subsection*{Numéricos (no notebook)}

\begin{enumerate}[label=\textbf{N\arabic*.},itemsep=4pt]
  \item Estenda o parser de METAR para rajadas (\texttt{G}), vento
        variável (\texttt{VRB}), \texttt{CAVOK} e visibilidade em
        milhas; teste com os 8 METARs fornecidos.
  \item Plote com o MetPy o modelo de estação dos METARs decodificados
        no N1, num mini-mapa do estado de São Paulo.
  \item No Caso~3, varra o raio de influência do Cressman e trace o
        erro RMS contra observações reservadas (validação cruzada);
        ache o raio ótimo.
  \item Reduza ao nível do mar as pressões da lista de estações
        fornecida (litoral a 800\,m) e trace as isóbaras antes e depois
        da redução — o mapa ``errado'' vs.\ o mapa sinótico.
\end{enumerate}

\vspace{1em}
\noindent\rule{\linewidth}{0.4pt}\\
{\scriptsize Material derivado de R.~Stull, \emph{Practical Meteorology}
(CC BY-NC-SA 4.0); este documento herda a mesma licença.}

\end{document}
